\documentclass[11pt,a4paper]{article}
\usepackage[margin=1in]{geometry}
\usepackage{libertine}
\usepackage[T1]{fontenc}
\usepackage[utf8]{inputenc}
\usepackage{amsfonts,amssymb,amsmath,amsthm}
% \usepackage[sort&compress]{natbib}
%\usepackage[english]{babel}
%\usepackage{eurosym}
\usepackage{calc}
\usepackage{tikz}
\usetikzlibrary{calc}

%\usepackage{verbatim,refcount}
%\usepackage{fancybox,soul,color}
\usepackage{url}
\usepackage{hyperref}
%\usepackage{subfig}
\usepackage{enumitem}
%\usepackage{epsfig}
%\usepackage[active]{srcltx}
%\usepackage[mathlines]{lineno}
\usepackage{xspace}
%\usepackage{multirow,tabu}
%\usepackage{slashed}
\usepackage{todonotes}
%\usepackage{eucal}
%\usepackage{mathtools}
\usepackage{thmtools}
\usepackage{mathrsfs}  
\usepackage{textpos}
\usepackage[normalem]{ulem}
\usepackage[capitalize, nameinlink]{cleveref}

%\usetikzlibrary{arrows,shapes,automata,backgrounds,petri, fit, decorations.pathreplacing,calligraphy}
%\DeclareMathOperator{\mad}{mad}
%\DeclareMathOperator{\ad}{ad}
%\DeclareMathOperator{\ntn}{ntn}
%\DeclareMathOperator{\func}{fun}
%\DeclareMathOperator{\degen}{degeneracy}
%\DeclareMathOperator{\avgd}{avgdeg}
%\DeclareMathOperator{\cwd}{cw}
%\DeclareMathOperator{\twin}{tww}
%\DeclareMathOperator{\fgn}{fgn}
\DeclareMathOperator{\VC}{2VCdim}
\DeclareMathOperator{\tw}{tw}
\DeclareMathOperator{\bsn}{bsn}
\DeclareMathOperator{\tpw}{tpw}
\DeclareMathOperator{\td}{td}
\DeclareMathOperator{\cc}{cc}
\DeclareMathOperator{\ddim}{ddim}

\newtheorem{problem}{Problem}[section]
\newtheorem{theorem}{Theorem}[section]
\newtheorem{lemma}[theorem]{Lemma}
\newtheorem{claim}[theorem]{Claim}
\newtheorem{corollary}[theorem]{Corollary}
\newtheorem{conjecture}[theorem]{Conjecture}
\theoremstyle{definition}
\newtheorem{definition}[theorem]{Definition}
\newcommand{\Oh}{\mathcal{O}}
\renewcommand{\varnothing}{\emptyset}

\def\Cc{\mathscr{C}}

\def\R{\mathbb{R}}
\def\C{\mathcal{C}}
\def\H{\mathcal{H}}
\def\P{\mathcal{P}}
\def\F{\mathcal{F}}
\def\N{\mathbb{N}}
\def\T{\mathcal{T}}
\def\D{\mathcal{D}}
\def\B{\mathrm{Ball}}
\def\Sol{\texttt{Sol}}
\newcommand{\Co}{\mathsf{Cont}}
\newcommand{\CMSO}{\textsf{CMSO}$_2$\xspace}
\newcommand{\cmsotwo}{\CMSO}
\newcommand{\VV}{$\bigvee\mkern-15mu\bigvee$}
\newcommand{\weight}{\mathsf{w}}

\newcommand{\Aa}{\mathcal{A}}
\newcommand{\Tt}{\mathcal{T}}
\newcommand{\Bb}{\mathcal{B}}
\newcommand{\Ss}{\mathcal{S}}
\newcommand{\Dd}{\mathcal{D}}
\newcommand{\Ww}{\mathcal{W}}

\newcommand{\wh}[1]{\widehat{#1}}

\newcommand{\bag}{\mathsf{bag}}
\newcommand{\dist}{\mathrm{dist}}
\newcommand{\rad}{\mathrm{rad}}

\newcommand{\Jadwiga}[2][]{\todo[color=yellow!80,inline,#1]{{\textbf{J:} #2}}}
\newcommand{\kk}[2][]{\todo[color=green!80,inline,#1]{{\textbf{kk:} #2}}}
\newcommand{\prz}[2][]{\todo[color=olive!80,#1]{{\textbf{prz:} #2}}}
\newcommand{\przin}[2][]{\todo[color=olive!80,#1,inline]{{\textbf{prz:} #2}}}
\newcommand{\mipilin}[2][]{\todo[color=blue!30,inline,#1]{{\textbf{MiPi:} #2}}}
\newcommand{\mipil}[2][]{\todo[color=blue!30,#1]{{\textbf{MiPi:} #2}}}
\newcommand{\mpil}[2][]{\todo[color=magenta!80,#1]{{\textbf{MP:} #2}}}

\newcommand{\tara}[2][]{\todo[color=violet!30,inline,#1]{{\textbf{Tara:} #2}}}

\newcommand{\ceil}[1]{\left\lceil #1 \right\rceil }

\renewcommand{\leq}{\leqslant}
\renewcommand{\geq}{\geqslant}
\renewcommand{\le}{\leqslant}
\renewcommand{\ge}{\geqslant}
\renewcommand{\mid}{~|~}

%\linenumbers

%%%%%%%%%%%%%%%%%%%%%%%%%%%%

%%%%%%%%%%%%%%%%%%%%%%%%%%%%%%%%%%%%%%%%%%%%%%%%%%%

\usepackage{framed}
\usepackage{tabularx}

\newlength{\RoundedBoxWidth}
\newsavebox{\GrayRoundedBox}
\newenvironment{GrayBox}[1]%
   {\setlength{\RoundedBoxWidth}{.93\columnwidth}
    \def\boxheading{#1}
    \begin{lrbox}{\GrayRoundedBox}
       \begin{minipage}{\RoundedBoxWidth}}%
   {   \end{minipage}
    \end{lrbox}
    \begin{center}
    \begin{tikzpicture}%
       \node(Text)[draw=black!20,fill=white,rounded corners,inner sep=2ex,text width=\RoundedBoxWidth]
             {\usebox{\GrayRoundedBox}};
        \coordinate(x) at (current bounding box.north west);
        \node [draw=white,rectangle,inner sep=3pt,anchor=north west,fill=white]
        at ($(x)+(6pt,.75em)$) {\boxheading};
    \end{tikzpicture}
    \end{center}}

\newenvironment{defproblemx}[1]{\noindent\ignorespaces%
                                \FrameSep=6pt%
                                \parindent=0pt%
%                 \vspace*{-1em}
                \begin{GrayBox}{#1}%
                \begin{tabular*}{\columnwidth}{!{\extracolsep{\fill}}@{\hspace{.1em}} >{\itshape} p{1.5cm} p{0.86\columnwidth} @{}}%
            }{
                \end{tabular*}%
                \end{GrayBox}%
                \ignorespacesafterend
%                 \vspace*{-1em}
            }

\newcommand{\problemTask}[3]{%
  \begin{defproblemx}{#1}
    Input: & #2 \\
    Task: & #3
  \end{defproblemx}
}

\newcommand{\parProblemTask}[4]{%
  \begin{defproblemx}{#1}
    Input:  & #2 \\
    Par.: & #3 \\
    Task: & #4
  \end{defproblemx}
}

\newcommand{\problemQuestion}[3]{%
  \begin{defproblemx}{#1}
    Input: & #2 \\
    Question: & #3
  \end{defproblemx}
}

\newcommand{\parProblemQuestion}[4]{%
  \begin{defproblemx}{#1}
    Input:  & #2 \\
    Par.: & #3 \\
    Question: & #4
  \end{defproblemx}
}

%\newcommand{\tara}[1]{{\color{violet}{(T: #1)}}}
\usepackage{mathtools}

\newcommand{\dd}{\hbox{-}}


\newenvironment{claimproof}[1][Proof of Claim.]{%
  \renewcommand{\qedsymbol}{$\lrcorner$}%
  \begin{proof}[#1]%
}{%
  \end{proof}%
}


\title{On coarse tree decompositions and coarse balanced separators}

\author{Tara Abrishami\thanks{University of Hamburg, Germany (\texttt{tara.abrishami@uni-hamburg.de}). Supported by the National Science Foundation Award Number DMS-2303251 and the Alexander von Humboldt Foundation.}
\and Jadwiga Czyżewska\thanks{University of Warsaw, Poland (\texttt{j.czyzewska@mimuw.edu.pl}).
Supported by Polish National Science Centre SONATA BIS-12 grant number 2022/46/E/ST6/00143.}
\and Kacper Kluk\thanks{University of Warsaw, Poland (\texttt{k.kluk@mimuw.edu.pl}).
Supported by Polish National Science Centre SONATA BIS-12 grant number 2022/46/E/ST6/00143.}
\and Marcin Pilipczuk\thanks{University of Warsaw, Poland (\texttt{m.pilipczuk@mimuw.edu.pl}).
Supported by Polish National Science Centre SONATA BIS-12 grant number 2022/46/E/ST6/00143.}
\and Michał Pilipczuk\thanks{University of Warsaw, Poland  (\texttt{michal.pilipczuk@uw.edu.pl}). Supported by the European Research Council (ERC) under the European Union’s Horizon 2020 research and innovation programme, grant agreement no 948057~(BOBR).}
\and Paweł Rzążewski\thanks{Warsaw University of Technology \& University of Warsaw, Poland (\texttt{pawel.rzazewski@pw.edu.pl}).}}

% Git access test :)

\begin{document}
\begin{titlepage}
\date{}
\maketitle

\begin{abstract}
It is known that there is a linear dependence between the treewidth of a graph and its {\em{balanced separator number}}: the smallest integer $k$ such that for every weighing of the vertices, the graph admits a balanced separator of size at most $k$. We investigate whether this connection can be lifted to the setting of coarse graph theory, where both the bags of the considered tree decompositions and the considered separators should be coverable by a bounded number of bounded-radius balls.

As the first result, we prove that if an $n$-vertex graph $G$ admits balanced separators coverable by $k$ balls of radius $r$, then $G$ also admits tree decompositions ${\cal T}_1$ and ${\cal T}_2$ such that:
\begin{itemize}[nosep]
    \item in ${\cal T}_1$, every bag can be covered by $\Oh(k\log n)$ balls of radius $r$; and
    \item in ${\cal T}_2$, every bag can be covered by $\Oh(k^2\log k)$ balls of radius $r(\log k+\log\log n+\Oh(1))$.
\end{itemize}
As the second result, we show that if we additionally assume that $G$ has doubling dimension at most $m$, then the functional equivalence between the existence of small balanced separators and of tree decompositions of small width can be fully lifted to the coarse setting. Precisely, we prove that for a positive integer $r$ and a graph $G$ of doubling dimension at most $m$, the following conditions are equivalent, with constants $k_1,k_2,k_3,k_4,\Delta_3,\Delta_4$ depending on each other and on $m$:
\begin{itemize}[nosep]
    \item $G$ admits balanced separators consisting of $k_1$ balls of radius $r$;
    \item $G$ has a tree decomposition with bags coverable by $k_2$ balls of radius $r$;
    \item $G$ has a tree-partition of maximum degree $\leq \Delta_3$ with bags coverable by $k_3$ balls of radius~$r$;
    \item $G$ is quasi-isometric to a graph of maximum degree $\leq \Delta_4$ and tree-partition~width $\leq k_4$.
\end{itemize}
\end{abstract}

\def\thepage{}
\thispagestyle{empty}
 \begin{textblock}{20}(-1.9, 4.6)
  \includegraphics[width=40px]{logo-erc}%
 \end{textblock}
 \begin{textblock}{20}(-1.9, 5.5)
  \includegraphics[width=40px]{logo-eu}%
 \end{textblock}
\end{titlepage}

\section{Introduction}


\begin{figure}[t]
\centering
\includegraphics[width=0.6\columnwidth]{figures/evaluation_desiderata_V5.pdf}
\vspace{-0.5cm}
\caption{\systemName is a platform for conducting realistic evaluations of code LLMs, collecting human preferences of coding models with real users, real tasks, and in realistic environments, aimed at addressing the limitations of existing evaluations.
}
\label{fig:motivation}
\end{figure}

\begin{figure*}[t]
\centering
\includegraphics[width=\textwidth]{figures/system_design_v2.png}
\caption{We introduce \systemName, a VSCode extension to collect human preferences of code directly in a developer's IDE. \systemName enables developers to use code completions from various models. The system comprises a) the interface in the user's IDE which presents paired completions to users (left), b) a sampling strategy that picks model pairs to reduce latency (right, top), and c) a prompting scheme that allows diverse LLMs to perform code completions with high fidelity.
Users can select between the top completion (green box) using \texttt{tab} or the bottom completion (blue box) using \texttt{shift+tab}.}
\label{fig:overview}
\end{figure*}

As model capabilities improve, large language models (LLMs) are increasingly integrated into user environments and workflows.
For example, software developers code with AI in integrated developer environments (IDEs)~\citep{peng2023impact}, doctors rely on notes generated through ambient listening~\citep{oberst2024science}, and lawyers consider case evidence identified by electronic discovery systems~\citep{yang2024beyond}.
Increasing deployment of models in productivity tools demands evaluation that more closely reflects real-world circumstances~\citep{hutchinson2022evaluation, saxon2024benchmarks, kapoor2024ai}.
While newer benchmarks and live platforms incorporate human feedback to capture real-world usage, they almost exclusively focus on evaluating LLMs in chat conversations~\citep{zheng2023judging,dubois2023alpacafarm,chiang2024chatbot, kirk2024the}.
Model evaluation must move beyond chat-based interactions and into specialized user environments.



 

In this work, we focus on evaluating LLM-based coding assistants. 
Despite the popularity of these tools---millions of developers use Github Copilot~\citep{Copilot}---existing
evaluations of the coding capabilities of new models exhibit multiple limitations (Figure~\ref{fig:motivation}, bottom).
Traditional ML benchmarks evaluate LLM capabilities by measuring how well a model can complete static, interview-style coding tasks~\citep{chen2021evaluating,austin2021program,jain2024livecodebench, white2024livebench} and lack \emph{real users}. 
User studies recruit real users to evaluate the effectiveness of LLMs as coding assistants, but are often limited to simple programming tasks as opposed to \emph{real tasks}~\citep{vaithilingam2022expectation,ross2023programmer, mozannar2024realhumaneval}.
Recent efforts to collect human feedback such as Chatbot Arena~\citep{chiang2024chatbot} are still removed from a \emph{realistic environment}, resulting in users and data that deviate from typical software development processes.
We introduce \systemName to address these limitations (Figure~\ref{fig:motivation}, top), and we describe our three main contributions below.


\textbf{We deploy \systemName in-the-wild to collect human preferences on code.} 
\systemName is a Visual Studio Code extension, collecting preferences directly in a developer's IDE within their actual workflow (Figure~\ref{fig:overview}).
\systemName provides developers with code completions, akin to the type of support provided by Github Copilot~\citep{Copilot}. 
Over the past 3 months, \systemName has served over~\completions suggestions from 10 state-of-the-art LLMs, 
gathering \sampleCount~votes from \userCount~users.
To collect user preferences,
\systemName presents a novel interface that shows users paired code completions from two different LLMs, which are determined based on a sampling strategy that aims to 
mitigate latency while preserving coverage across model comparisons.
Additionally, we devise a prompting scheme that allows a diverse set of models to perform code completions with high fidelity.
See Section~\ref{sec:system} and Section~\ref{sec:deployment} for details about system design and deployment respectively.



\textbf{We construct a leaderboard of user preferences and find notable differences from existing static benchmarks and human preference leaderboards.}
In general, we observe that smaller models seem to overperform in static benchmarks compared to our leaderboard, while performance among larger models is mixed (Section~\ref{sec:leaderboard_calculation}).
We attribute these differences to the fact that \systemName is exposed to users and tasks that differ drastically from code evaluations in the past. 
Our data spans 103 programming languages and 24 natural languages as well as a variety of real-world applications and code structures, while static benchmarks tend to focus on a specific programming and natural language and task (e.g. coding competition problems).
Additionally, while all of \systemName interactions contain code contexts and the majority involve infilling tasks, a much smaller fraction of Chatbot Arena's coding tasks contain code context, with infilling tasks appearing even more rarely. 
We analyze our data in depth in Section~\ref{subsec:comparison}.



\textbf{We derive new insights into user preferences of code by analyzing \systemName's diverse and distinct data distribution.}
We compare user preferences across different stratifications of input data (e.g., common versus rare languages) and observe which affect observed preferences most (Section~\ref{sec:analysis}).
For example, while user preferences stay relatively consistent across various programming languages, they differ drastically between different task categories (e.g. frontend/backend versus algorithm design).
We also observe variations in user preference due to different features related to code structure 
(e.g., context length and completion patterns).
We open-source \systemName and release a curated subset of code contexts.
Altogether, our results highlight the necessity of model evaluation in realistic and domain-specific settings.





% !TEX root =  ../main.tex
\section{Background on causality and abstraction}\label{sec:preliminaries}

This section provides the notation and key concepts related to causal modeling and abstraction theory.

\spara{Notation.} The set of integers from $1$ to $n$ is $[n]$.
The vectors of zeros and ones of size $n$ are $\zeros_n$ and $\ones_n$.
The identity matrix of size $n \times n$ is $\identity_n$. The Frobenius norm is $\frob{\mathbf{A}}$.
The set of positive definite matrices over $\reall^{n\times n}$ is $\pd^n$. The Hadamard product is $\odot$.
Function composition is $\circ$.
The domain of a function is $\dom{\cdot}$ and its kernel $\ker$.
Let $\mathcal{M}(\mathcal{X}^n)$ be the set of Borel measures over $\mathcal{X}^n \subseteq \reall^n$. Given a measure $\mu^n \in \mathcal{M}(\mathcal{X}^n)$ and a measurable map $\varphi^{\V}$, $\mathcal{X}^n \ni \mathbf{x} \overset{\varphi^{\V}}{\longmapsto} \V^\top \mathbf{x} \in \mathcal{X}^m$, we denote by $\varphi^{\V}_{\#}(\mu^n) \coloneqq \mu^n(\varphi^{\V^{-1}}(\mathbf{x}))$ the pushforward measure $\mu^m \in \mathcal{M}(\mathcal{X}^m)$. 


We now present the standard definition of SCM.

\begin{definition}[SCM, \citealp{pearl2009causality}]\label{def:SCM}
A (Markovian) structural causal model (SCM) $\scm^n$ is a tuple $\langle \myendogenous, \myexogenous, \myfunctional, \zeta^\myexogenous \rangle$, where \emph{(i)} $\myendogenous = \{X_1, \ldots, X_n\}$ is a set of $n$ endogenous random variables; \emph{(ii)} $\myexogenous =\{Z_1,\ldots,Z_n\}$ is a set of $n$ exogenous variables; \emph{(iii)} $\myfunctional$ is a set of $n$ functional assignments such that $X_i=f_i(\parents_i, Z_i)$, $\forall \; i \in [n]$, with $ \parents_i \subseteq \myendogenous \setminus \{ X_i\}$; \emph{(iv)} $\zeta^\myexogenous$ is a product probability measure over independent exogenous variables $\zeta^\myexogenous=\prod_{i \in [n]} \zeta^i$, where $\zeta^i=P(Z_i)$. 
\end{definition}
A Markovian SCM induces a directed acyclic graph (DAG) $\mathcal{G}_{\scm^n}$ where the nodes represent the variables $\myendogenous$ and the edges are determined by the structural functions $\myfunctional$; $ \parents_i$ constitutes then the parent set for $X_i$. Furthermore, we can recursively rewrite the set of structural function $\myfunctional$ as a set of mixing functions $\mymixing$ dependent only on the exogenous variables (cf. \cref{app:CA}). A key feature for studying causality is the possibility of defining interventions on the model:
\begin{definition}[Hard intervention, \citealp{pearl2009causality}]\label{def:intervention}
Given SCM $\scm^n = \langle \myendogenous, \myexogenous, \myfunctional, \zeta^\myexogenous \rangle$, a (hard) intervention $\iota = \operatorname{do}(\myendogenous^{\iota} = \mathbf{x}^{\iota})$, $\myendogenous^{\iota}\subseteq \myendogenous$,
is an operator that generates a new post-intervention SCM $\scm^n_\iota = \langle \myendogenous, \myexogenous, \myfunctional_\iota, \zeta^\myexogenous \rangle$ by replacing each function $f_i$ for $X_i\in\myendogenous^{\iota}$ with the constant $x_i^\iota\in \mathbf{x}^\iota$. 
Graphically, an intervention mutilates $\mathcal{G}_{\mathsf{M}^n}$ by removing all the incoming edges of the variables in $\myendogenous^{\iota}$.
\end{definition}

Given multiple SCMs describing the same system at different levels of granularity, CA provides the definition of an $\alpha$-abstraction map to relate these SCMs:
\begin{definition}[$\abst$-abstraction, \citealp{rischel2020category}]\label{def:abstraction}
Given low-level $\mathsf{M}^\ell$ and high-level $\mathsf{M}^h$ SCMs, an $\abst$-abstraction is a triple $\abst = \langle \Rset, \amap, \alphamap{} \rangle$, where \emph{(i)} $\Rset \subseteq \datalow$ is a subset of relevant variables in $\mathsf{M}^\ell$; \emph{(ii)} $\amap: \Rset \rightarrow \datahigh$ is a surjective function between the relevant variables of $\mathsf{M}^\ell$ and the endogenous variables of $\mathsf{M}^h$; \emph{(iii)} $\alphamap{}: \dom{\Rset} \rightarrow \dom{\datahigh}$ is a modular function $\alphamap{} = \bigotimes_{i\in[n]} \alphamap{X^h_i}$ made up by surjective functions $\alphamap{X^h_i}: \dom{\amap^{-1}(X^h_i)} \rightarrow \dom{X^h_i}$ from the outcome of low-level variables $\amap^{-1}(X^h_i) \in \datalow$ onto outcomes of the high-level variables $X^h_i \in \datahigh$.
\end{definition}
Notice that an $\abst$-abstraction simultaneously maps variables via the function $\amap$ and values through the function $\alphamap{}$. The definition itself does not place any constraint on these functions, although a common requirement in the literature is for the abstraction to satisfy \emph{interventional consistency} \cite{rubenstein2017causal,rischel2020category,beckers2019abstracting}. An important class of such well-behaved abstractions is \emph{constructive linear abstraction}, for which the following properties hold. By constructivity, \emph{(i)} $\abst$ is interventionally consistent; \emph{(ii)} all low-level variables are relevant $\Rset=\datalow$; \emph{(iii)} in addition to the map $\alphamap{}$ between endogenous variables, there exists a map ${\alphamap{}}_U$ between exogenous variables satisfying interventional consistency \cite{beckers2019abstracting,schooltink2024aligning}. By linearity, $\alphamap{} = \V^\top \in \reall^{h \times \ell}$ \cite{massidda2024learningcausalabstractionslinear}. \cref{app:CA} provides formal definitions for interventional consistency, linear and constructive abstraction.
\section{Distance graphs}\label{sec:quasi-isometries}
 In this section we describe a generic construction of a quasi-isometry between a graph and its ``coarsening'' with respect to some magnitude of distances. The construction can be considered folklore, see e.g.~\cite[Observation 2.1]{coarse2023} and further references mentioned there, but as we will later use its specific properties, we describe it in details. %We also choose to describe it in the setting of metric spaces for the sake of generality.

 Let $G$ be a graph and $r\in \N_{>0}$. Suppose $I$ is an inclusion-maximal distance-$r$ independent set in~$G$. The {\em{$(I,r)$-distance graph}} of $G$ is the edge-weighted graph $H=H(G,I,r)$ defined as follows:
 \begin{itemize}[nosep]
 \item the vertex set of $H$ is $I$; and
 \item for every two distinct vertices $u,v\in I$ satisfying $\dist(u,v)\leq 3r$, in $H$ we add an edge $uv$ of weight $3r$.
 \end{itemize}
 Thus, all edges of $H$ have weight $3r$. %Note that a priori $H$ may be an infinite graph in case $I$ is infinite, but we will apply the construction only in the context of $(X,d)$ being the distance metric of a finite~graph.
 
 %, we define a graph $H$ with vertex set $I$ and with an edge between two vertices $u, v$ of $I$ if $d(u, v) \leq 3r$. 
 %Each edge of $H$ has weight $3r$.
 %Let $d_H$ be a shortest path metric on graph $H$. Then $(V(H), d_H)$ is a metric space. We will say that $H$ and $(V(H), d_H)$ are \emph{derived} from the tuple $(X,d,I,r)$. In the following we will study the properties of $(V(H), d_H)$. 

%In the following we will slightly abuse the notation, sometimes referring to vertices of $H$ as to elements of $X$. 
%We will also consider the  distance between them according to the metric $d$ or $d_H$.

Let us first note that any distance graph derived from a graph of bounded doubling dimension has bounded maximum degree.

\begin{lemma}\label{clm:bdddim_implies_bddegree}
    Let $G$ be a graph of doubling dimension $m$, $I$ be an inclusion-wise maximal distance-$r$ independent set in $G$ for some $r\in \N_{>0}$, and $H$ be the $(I,r)$-distance graph of $G$. Then $\Delta(H)<2^{3m}$.
\end{lemma}
\begin{proof}
    Consider a vertex $u \in I$ and let $N_H[u]$ be the closed neighborhood of $u$ in $H$, i.e., the set comprising $u$ and all its neighbors. Since two vertices $x,y \in I$ are adjacent in $H$ only if $\dist_G(x,y) \leq 3r$, it follows that $N_H[u] \subseteq \B_G(u, 3r)$. Since $G$ has doubling dimension at most $m$, there are $2^{3m}$ balls $B_1, \hdots, B_{2^{3m}}$ of $X$, each of radius $3r/8$, such that $\B_G(u, 3r) \subseteq \bigcup_{i=1}^{2^{3m}} B_i$. Observe that for each $B_i$, the maximum distance between any two vertices in $B_i$ is at most $6r/8$. Since the vertex set of $H$ is a distance-$r$ independent set in $G$, we have $\dist_G(x,y) > r$ for any distinct $x,y \in N_H[u]$, so each ball $B_i$ contains at most one vertex of $N_H[u]$. As their union contains every vertex of $N_H[u]$, we conclude that $|N_H[u]| \leq 2^{3m}$, hence the degree of $u$ is strictly smaller than $2^{3m}$.
\end{proof}

We now show that every graph is suitably quasi-isometric to any its distance graph.

\begin{lemma}\label{lem:quasi_isometry_to_derived_graph}
    Suppose that $G$ is a graph, $I$ is an inclusion-wise maximal distance-$r$ independent set in $G$ for some $r\in \N_{>0}$, and $H$ is the $(I,r)$-distance graph of $G$. For every $u\in V(G)$, let $\varphi(u)$ be an arbitrary vertex of $I$ such that $\dist_G(u,\varphi(u))\leq r$ (such a vertex exists by the maximality of $I$). Then $\varphi$ is a $(3,3r)$-quasi-isometry from $G$ to $H$.
\end{lemma}
\begin{proof}
    For the second condition of quasi-isometry, observe that for every $w\in V(H)=I$, we have $\dist_G(w,\varphi(w))\leq 3r$. So it remains to  show that the first property is also satisfied. We may assume without loss of generality that $G$ is connected.

    %Let us fix an arbitrary order on $I$ and consider the Voronoi partition of $X$ with respect to the set $I$. Define a mapping $\varphi$ such that every point $x$ of $X$ is mapped to the unique vertex $f$ of $I$ such that $x \in R_V(f)$. Now $\varphi$ is indeed a mapping $\varphi: X \mapsto V(H)$. We show that $\varphi$ is a $(3, 3r)$-quasi-isometry. Since $f \in R_V(f)$ for all $f \in V(H)$, the mapping $\varphi$ is surjective and therefore satisfies the second property. We next show that the first property is also satisfied.


    % \tara{Since $I$ is an inclusion-maximal dist-$r$ independent set of $X$, it follows that $d(v, \varphi(v)) \leq r$ for every $v \in X$. Fix $u, v \in X$. Let $P = d_X(u, v)$ and $Q = d_Y(\varphi(u), \varphi(v))$. Let $q_1 \dd q_2 \dd \hdots \dd q_Q$ be a shortest path in $H$ from $q_1 = \varphi(u)$ to $q_Q = \varphi(v)$. Since two vertices are adjacent in $H$ only if their distance is at most $3r$ in $X$, it follows that there is a curve of length at most $r + 3r \cdot Q + r$ from $u$ to $v$ in $X$. Therefore, $P/3r - 2/3 \leq Q$. %$P \leq 3r\cdot Q + 2r$. 
    % %Next, let $p_1 \dd \hdots \dd p_S$ be the image of a shortest curve in $X$ from $u$ to $v$ under $\varphi$, so $p_1 = \varphi(u)$ and $p_S = \varphi(v)$. 
    % }

    %each edge of H has weight 3r and this is exactly the reason - we want to get rid of r here in the multiplicative factor

    
    %Since $I$ is an inclusion-maximal distance-$r$ independent set, every element of $X$ is at distance at most $r$ from an element of $I$, so $d(v,\varphi(v))\leq r$ for all $x \in X$.

    Consider any $u,v\in V(G)$ and let $P$ be a shortest path connecting $u$ and $v$. Along $P$ we may choose vertices $x_0, x_1, \dots, x_k, x_{k+1}$ in order so that $x_0=u$, $x_{k+1}=v$ and
    \begin{itemize}[nosep]
        \item for each $i\in \{0,1\dots, k-1\}$ we have $\dist_G(x_i, x_{i+1})=r$, and
        \item $\dist_G(x_k, x_{k+1})\leq r$.
    \end{itemize}
    Note that $\sum_{i=0}^k\ \dist_G(x_i, x_{i+1}) = \dist_G(u,v)$ and we have $kr<\dist_G(u,v) \leq (k+1)r$.

    For each $i\in\{0,1,\dots, k+1\}$, let $v_i=\varphi(x_i)\in I$. Note that $\dist_G(x_i,v_i)\leq r$ by construction. 
    We claim that for every $i\in \{0,1,\ldots,k\}$, the vertices $v_i$ and $v_{i+1}$ are adjacent in the graph $H$. 
    Indeed, $\dist_G(v_i, v_{i+1})\leq \dist_G(v_i, x_i) + \dist_G(x_i, x_{i+1}) + \dist_G(x_{i+1}, v_{i+1}) \leq 3r$ by the triangle inequality, so we have $v_iv_{i+1}\in E(H)$. 
    Consequently, $v_0,\ldots,v_{k+1}$ is a walk that connects $\phi(u) = v_0$ and $\phi(v)=v_{k+1}$ in $H$.
    Recalling that each edge in $H$ has weight $3r$, we may now estimate $\dist_H(\varphi(u), \varphi(v))$ as follows:
    %\[\dist_H(\varphi(u), \varphi(v))=\dist_H(x_0, x_{k+1})\leq \sum_{i=0}^k\ \dist_H(x_i, x_{i+1}) \leq 3r\cdot(k+1)\leq 3\cdot \dist_G(u, v) +3r.\]
    \[\dist_H(\varphi(u), \varphi(v))=\dist_H(v_0, v_{k+1})\leq \sum_{i=0}^k\ \dist_H(v_i, v_{i+1}) \leq 3r\cdot(k+1)\leq 3\cdot \dist_G(u, v) +3r.\]
    
    Next, since $\dist_G(x, y) \leq 3r$ for each edge $xy \in E(H)$ and every edge of $H$ has weight $3r$, it follows that $\dist_G(x,y) \leq \dist_H(x,y)$ for all $x,y \in I$. Now, we can estimate $\dist_G(u,v)$:
    \[\dist_G(u,v) \leq \dist_G(u, \varphi(u)) + \dist_G(\varphi(u), \varphi(v)) + \dist_G(\varphi(v), u) \leq r + \dist_H(\varphi(u), \varphi(v)) + r.\]
    Thus, we get the following inequalities:
    \[\dist_G(u,v) - 2r \leq \dist_H(\varphi(u), \varphi(v)) \leq 3\cdot \dist_G(u, v) + 3r,\]
    which means that also the following is true
    \[\tfrac{1}{3}\cdot \dist_G(u,v) - 3r \leq \dist_H(\varphi(u), \varphi(v)) \leq 3\cdot \dist_G(u, v) + 3r.\]
    All in all, we conclude that $\varphi$ is indeed a $(3,3r)$-quasi-isometry.
\end{proof}

  
 \section{Bounded doubling dimension}

In this section we prove \Cref{thm:equivalences}. We start with the following lemma that encapsulates the key observation: in the presence of a bound on the doubling dimension, the distance-$r$ balanced separators in $G$ can be translated into balanced separators in the corresponding distance graph.

 \begin{lemma}\label{lem:bdtw}
     Let $G$ be a graph with doubling dimension at most $m$, for some $m\in \N$. Suppose that for some $r\in \N_{>0}$, the graph $G$ has distance-$r$ balanced separator number at most $k$. Then, for every inclusion-wise maximal distance-$r$ independent set $I$ in $G$, the distance graph $H(G,I,r)$ has treewidth at most  $3k\cdot 2^{6m}$.
 \end{lemma}
\begin{proof}
    Denote $H=H(G,I,r)$ for brevity. By \cref{lem:bsn}, it suffices to prove that the balanced separator number of $H$ is bounded by $k\cdot 2^{6m}$. 
    Consider any weight function $\mu_H\colon V(H)\to \R_{\geq 0}$ on $H$. Let $\mu_G\colon V(G) \to \R_{\geq 0}$ be the weight function on $G$ defined by setting $\mu_G(u)=\mu_H(u)$ for all $u\in I$, and $\mu_G(u)=0$ for all $u\notin I$. Since $G$ has distance-$r$ balanced separator number at most $k$, there exist radius-$r$ balls $B_1,B_2,\dots, B_{k}$ in $G$ such that denoting $X\coloneqq \bigcup_{i=1}^{k} B_i$, every connected component of $G-X$ has weight at most~$\frac{1}{2}\mu_G(V(G))$.

    Let $u_1, u_2, \dots, u_{k}$ be the centers of balls $B_1,B_2,\ldots,B_k$, respectively. Since $I$ is an inclusion-wise maximal distance-$r$ independent set, we can find vertices $v_1,\ldots,v_r\in I$ such that $\dist_G(u_i,v_i)\leq r$, for each $i\in \{1,\ldots,k\}$. Further, let $N_H^2[v_i]$ be the second neighborhood of $v_i$ in $H$: the set consisting of $v_i$, the neighbors of $v_i$, and the neighbors of those neighbors.
    We define
    \[X_H\coloneqq \bigcup_{i=1}^k N_H^2[v_i]\subseteq V(H)=I.\]
    By \cref{clm:bdddim_implies_bddegree}, $\Delta(H)<2^{3m}$. It follows that \[|N_H^2[v_i]|\leq 1+\Delta(H)+\Delta(H)\cdot (\Delta(H)-1)=1+(\Delta(H))^2\leq 2^{6m},\quad\textrm{for each }i\in \{1,\ldots,k\}.\]
    Therefore, we have $|X_H|\leq k\cdot 2^{6m}$.
    
    We will show that $X_H$ is a balanced separator for~$\mu_H$.
    Note that since $X$ is a balanced separator for $\mu_G$, it suffices to prove the following claim: every two vertices $u,v\in I\setminus X_H$ that are in different components of $G-X$, are also in different components of~$H-X_H$.

    Consider any path $Q$ connecting $u$ and $v$ in $H$.
    Let $P$ be a walk connecting $u$ and $v$ in $G$ obtained by replacing every edge of $Q$, say $xy$, with a path $R_{xy}$ in $G$ of length at most $3r$. Such a path exists by the construction of $H$ as it is the $(I,r)$-distance graph of $G$.
    Since $u$ and $v$ are in different components of $G-X$, the walk $P$ must intersect $X$. This means that for some edge $xy$ of $Q$ and $i\in \{1,\ldots,k\}$, there exists a vertex $w\in V(R_{xy})$ such that $\dist_G(w,u_i)\leq r$, implying that $\dist_G(w,v_i)\leq 2r$.
    Since the length of $R_{xy}$ is at most $3r$, we have $\dist_G(w,x)\leq \frac{3}{2}r$ or $\dist_G(w,y)\leq \frac{3}{2}r$, implying that $\dist_G(v_i,x)\leq \frac{7}{2}r$ or $\dist_G(v_i,y)\leq \frac{7}{2}r$. Without loss of generality assume the former. Then on a shortest path between $x$ and $v_i$ we may find a vertex $z$ such that $\dist_G(x,z)\leq 2r$ and $\dist_G(v_i,z)\leq 2r$. Since $I$ is an inclusion-wise maximal distance-$r$ independent set, there exists some $z'\in I$ such that $\dist_G(z,z')\leq r$. In particular we have $\dist_G(x,z')\leq 3r$ and $\dist(z',v_i)\leq 3r$, which means that $x,z'$ are equal or adjacent in $H$, and also $z',v_i$ are equal or adjacent in $H$. We conclude that $x\in N_H^2[v_i]$, so $x\in X_H$ and therefore $Q$ intersects $X_H$. As $Q$ was chosen arbitrarily, we conclude that $u$ and $v$ lie in different components of $H-X_H$ and we are~done. 
\end{proof}
    
    
    
%    By triangle inequality we have $B_i\subseteq \B_G(v_i,2r)$, hence
%    \[\bigcup_{i=1}^{k} B_i \subseteq \bigcup_{i=1}^{k} \B_{G}(v_i, 2r).\] 
%    Since So each component of $G - \bigcup_{i=1}^{k_4} \B_{(X,d)}(v_i, 2r)$ has weight at most $\frac{\mu(X)}{2}$. Therefore each component of $H-\bigcup_{i=1}^{k_4}\B_{(V(H), d_H)}(v_i, 1)$ has weight at most $\frac{\mu(V(H))}{2}$ (recall that $z_1z_2\notin E(H)$ implies $d(z_1, z_2) > 3r$). Thus $\bigcup_{i=1}^{k_4} \B_{(V(H, d_H)}(v_i, 1)$ is a balanced separator for $\mu_H$ of $H$ of size $k_4 \cdot 2^m$, as each vertex of $H$ has degree bounded by $2^m$ by Claim~\ref{clm:bdddim_implies_bddegree}. 

%    We showed that for each weight function $\mu_H\colon V(H)\to \mathbb{R}_{\geq 0}$ there exists a balanced separator of $H$ of size $2^mk_4$. Thus $H$ has treewidth equal $O(k_42^m)$.



%\begin{theorem}\label{thm:equivalences}
%    Let $\Cc$ be a class of graphs of doubling dimension bounded by $m$, for some $m\in \N$. Then the following conditions are equivalent for every $r\in \N_{>0}$:
%    \begin{enumerate}[label=(\arabic*),ref=(\arabic*)]
%        \item\label{pr:tpw} There exist $k_1,\Delta_1\in \N$ such that every member of $\Cc$ has a tree-partition of spread $r$, maximum degree at most $\Delta_1$, and with $(k_1,r)$-coverable bags.
%        \item\label{pr:tw} There exists $k_2\in \N$ such that every member of $\Cc$ has a tree decomposition with $(k_2,r)$-coverable~bags.
 %       \item\label{pr:bsn} There exists $k_3\in \N$ such that every member of $\Cc$ has distance-$r$ balanced separator number at most $k_3$.
 %       \item\label{pr:qi} There exist $k_4,\Delta_4\in \N$ such that every member of $\Cc$ is $(3,3r)$-quasi-isometric with an edge-weighted graph of tree-partition width at most $k_4$, maximum degree at most $\Delta_4$, and every edge of weight $3r$.
 %       \item\label{pr:qia} There exist $k_5,\Delta_5\in \N$ and $\alpha,\beta,\gamma\in \R_{>0}$ such that every member of $\Cc$ is $(\alpha, \beta r)$-quasi-isometric with an edge-weighted graph of tree-partition width at most $k_5$, maximum degree at most $\Delta_5$, and every edge of weight at least $\gamma r$.
 %   \end{enumerate}
%\end{theorem}
\thmequivalences*
\begin{proof}
    \noindent{\bf{\ref{pr:tpw}$\Rightarrow$ \ref{pr:tw}.}} We use the standard construction sketched in \cref{sec:prelim}. Consider any $G\in \Cc$ and let $(T,\bag)$ be a tree-partition of $G$ with $(k_1,r)$-coverable bags (we will not use the assumption about the spread or the maximum degree). We construct a new tree $T'$ by subdividing every edge of $T$, say $xy$, with a new vertex~$z_{xy}$. On $T'$ we define a new bag function $\bag'$ as follows: $\bag'(x)=\bag(x)$ for each $x\in V(T)$ and $\bag'(z_{xy})=\bag(x)\cup \bag(y)$ for each $xy\in E(T)$. It is straightforward to see that $(T',\bag')$ is a tree decomposition of $G$, and its bag are $(k_2,r)$-coverable, where $k_2\coloneqq 2k_1$.

    \bigskip

    \noindent{\bf{\ref{pr:tw}$\Rightarrow$ \ref{pr:bsn}.}}
     Again, we use the standard argument of finding a balanced bag in a tree decomposition. Consider a graph $G\in \Cc$ and let
     $(T,\bag)$ be a tree decomposition of $G$ with $(k_2,r)$-coverable bags.  
    Consider also any weight function $\mu\colon V(G)\to \R_{\geq 0}$. We orient the edges of $T$ according to $\mu$ as follows. Consider an edge $xy$ of $T$; then removal of $xy$ from $T$ disconnects $T$ into two subtrees, say $T_x$ containing $x$ and $T_y$ containing $y$. If $\mu(\bigcup_{x'\in V(T_x)} \bag(x)) < \mu(\bigcup_{y'\in T_y} \bag(y))$, then orient $xy$ towards~$y$, and if $\mu(\bigcup_{x'\in V(T_x)} \bag(x)) > \mu(\bigcup_{y'\in T_y} \bag(y))$, then orient $xy$ towards $x$; in case of a tie, orient $xy$ in any way.
    Since every orientation of a tree has a sink, there exists a node $x\in V(T)$ such that all edges incident to $x$ are oriented towards $x$.
    Then each component of $G-\bag(x)$ has $\mu$-weight at most~$\frac{1}{2}\mu(V(G))$, as otherwise the edge connecting $x$ with the tree of $T-x$ containing this component would need to be oriented outwards from $x$. Since $\bag(x)$ is $(k_2,r)$-coverable, we conclude that $\bag(x)$ is a $(k_2,r)$-coverable balanced separator for $\mu$. And as $\mu$ was chosen arbitrarily, $G$ has distance-$r$ separation number bounded by $k_3\coloneqq  k_2$.

    \bigskip

    \noindent{\bf{\ref{pr:bsn}$\Rightarrow$ \ref{pr:qi}.}}
    Consider a graph $G\in \Cc$.
    Let $I$ be an inclusion-maximal distance-$r$ independent set in $G$, and let $H\coloneqq H(G,I,r)$ be the $(I,r)$-distance graph of $G$. Then $G$ is $(3,3r)$-quasi-isometric with~$H$ (by \cref{lem:quasi_isometry_to_derived_graph}), and $H$ has maximum degree bounded by $\Delta_4\coloneqq 2^{3m}-1$ (by \cref{clm:bdddim_implies_bddegree}) and tree-partition-width bounded by $k_4\coloneqq \frac{35}{4}\Delta_4\cdot (3k_3\cdot 2^{6m}+1)$ (by \cref{thm:tw-tpw,lem:bdtw}).
    
    %there exists a $(3,3r)$-quasi-isometry $\varphi:X\to V(H)$ . We will show that there is a tree-partition of $H$ of width bounded by a function of $k_4$, $r$ and $m$.

    %By Lemma~\ref{lem:bdtw} $H$ has treewidth equal $O(2^mk)$.
    %Since $H$ has maximum degree bounded by $2^m$, there exists a tree-partition $\mathcal{T}=(T, \{\beta_x\}_{x\in V(T)})$ of $H$ of width $24(\tw(H)+1)\Delta(H)=O(k_4\cdot2^{m}) := k_1$ and maximum degree bounded by $\Delta_1 := 24(\tw(G) + 1)\Delta(G)^2 = O(k_4\cdot 2^{m+1})$.
%    \newline
%    \newline
%    $\mathbf{1\Rightarrow 2}$.
%    Let us just take $\alpha=\beta =3$, $k_2=k_1$ and $\Delta_2=\Delta_1$. Then the statement is clearly fulfilled.

    \bigskip

    \noindent{\bf{\ref{pr:qi}$\Rightarrow$ \ref{pr:qia}.}} Trivial, take $k_5\coloneqq k_4$, $\Delta_5\coloneqq \Delta_4$, $\alpha\coloneqq 3$, $\beta\coloneqq 3$, and $\gamma\coloneqq 3$.

    \bigskip
    
    \noindent{\bf{\ref{pr:qia}$\Rightarrow$ \ref{pr:tpw}.}} We may assume without loss of generality that $\alpha,\beta\geq 1$. Consider a graph $G\in \Cc$ and let 
    $H$ be an edge-weighted graph with $\tpw(H)\leq k_5$, $\Delta(H)\leq \Delta_5$, and every edge of weight at least $\gamma r$, such that there is an $(\alpha,\beta r)$-quasi-isometry $\varphi \colon V(G)  \to V(H)$.
    Let $\Tt=(T,\bag)$ be a tree-partition of $H$ of width at most $k_5$.
    Note that we may assume that $T$ has maximum degree at most $\Delta_5'\coloneqq k_5\cdot \Delta_5$, for the total number of neighbors of vertices contained in a single bag of $\Tt$ is at most $\Delta_5'$.

    Our goal is to construct a suitable tree-partition of $G$. We first define a new bag function $\bag'$ by naturally pulling $\bag$ through the quasi-isometry $\varphi$:
    \[\bag'(x)\coloneqq \{u\in V(G)\mid \varphi(u)\in \bag(x)\} \qquad \textrm{for each }x\in V(T).\]
        
    We claim that for each $x\in V(T)$, $\bag(x)$ is $(k_5',r)$-coverable, where $k_5'\coloneqq k_5\cdot 2^{m\lceil\log \alpha\beta\rceil}$. For this, observe that if we have any pair of vertices $u,v\in V(G)$ with $\varphi(u)=\varphi(v)$, then 
    \[\tfrac{1}{\alpha}\dist_G(u,v) - \beta r \leq \dist_H(\varphi(u), \varphi(v)) = 0,\qquad\textrm{implying}\qquad \dist_G(u, v)\leq \alpha\beta r.\]

    Therefore, for each $v\in \bag'(x)$ with non-empty $\varphi^{-1}(v)$ we can select any vertex $u_v\in \varphi^{-1}(v)$, and then $\varphi^{-1}(v)\subseteq \B_G(u_v, \alpha\beta r)$.
    Since $\bag(x)$ consists of at most $k_5$ elements, it follows that $\bag'(x)$ is $(k_5,\alpha\beta r)$-coverable. Since $G$ has doubling dimension at most $m$, each radius-$\alpha\beta r$ ball in $G$ can be covered by $2^{m\lceil\log \alpha\beta\rceil}$ radius-$r$ balls. It follows that $\bag'(x)$ is $(k_5',r)$-coverable, as claimed.

    It would be natural to consider $(T,\bag')$ as the sought tree-partition, but it is not clear that this is a tree-partition at all. Instead, we will construct a suitable ``coarsening'' of $(T,\bag')$ as follows; see also \cref{fig:tree}. Set
    \[p\coloneqq \lceil (\alpha+\beta)\gamma^{-1}\rceil.\]
    Root $T$ in an arbitrary node $z$ and let
    \[I\coloneqq \{y\in V(T)\mid \dist_T(y,z)\equiv 0 \bmod p\}.\]
    Here, the edges of $T$ are considered to be of unit length. Further, define
    \[C_z\coloneqq \{x\in V(T)\mid \dist_T(x,z)<2p\},\]
    and for each $y\in I\setminus \{z\}$, 
    \[C_y\coloneqq \{x\in V(T)\mid p\leq \dist_T(x,y)<2p\textrm{ and }\dist_T(x,z)>\dist_T(y,z)\}.\]
    Note that $\{C_y \mid  y \in I\}$ is a partition of $V(T)$. Since $T$ has maximum degree at most $\Delta_5'$, we have
    \begin{equation}\label{eq:wydra}|C_y|\leq 1+\Delta_5'+(\Delta_5')^2+\ldots+(\Delta_5')^{2p-1}\eqqcolon M\qquad\textrm{for each }y\in I.\end{equation}
    Next, define $T''$ to be the tree on the node set $I$ where $y,y'\in I$ are adjacent in $T''$ whenever \[\dist_T(y,y')=p\qquad\textrm{and}\qquad|\dist_T(y,z)-\dist_T(y',z)|=p.\]
    Then, from the construction we immediately obtain the following: 
    \begin{equation}\label{eq:bobr}
        \textrm{for all $x,x'\in V(T)$, if $x\in C_y$ and $x'\in C_{y'}$ with $y\neq y'$ and $yy'\notin E(T')$, then $\dist_T(x,x')>p$.}
    \end{equation}
    Note also that $T''$ has maximum degree bounded by $\Delta_1\coloneqq 1+(\Delta'_5)^p$.

\begin{figure}[t]
    \centering
    \includegraphics[scale=0.72, page=1]{tree.pdf}
    \includegraphics[scale=0.72, page=2]{tree.pdf}
    \caption{The construction of $T''$ in the proof of \cref{thm:equivalences}. \textbf{Left:} The definition of the sets $I$ (red), $C_z$, and $C_y$ for a vertex $y \in V(T)$ (both sets blue). \textbf{Right:} The vertices and edges of $T''$.}
    \label{fig:tree}
\end{figure}

    We endow $T''$ with bag function $\bag''$ defined as follows:
    \[\bag''(y)\coloneqq \bigcup_{x\in C_y} \bag'(x),\qquad\textrm{for all }y\in I=V(T'').\]
    Since $\bag'(x)$ is $(k_5',r)$-coverable for each $x\in V(T)$, from \eqref{eq:wydra} it follows that $\bag''(y)$ is $(k_1,r)$-coverable for each $y\in V(T)$, where $k_1\coloneqq M\cdot k_5'$.
    It remains to show that $\Tt''\coloneqq (T'',\bag'')$ is a tree-partition of $G$ of spread $r$.
    
    First, note that $\{\bag'(y) \mid y \in V(T'') \}$ is a partition of $V(G)$.
    Indeed, this follows from the fact that $\varphi(v) \in V(H)$ is  uniquely defined for every $v \in V(G)$ and $\{\bag(x) \mid x \in V(T) \}$ is a partition of $V(H)$.
    Now consider any $u,v\in V(G)$ with $u\in \bag''(y)$ and $v\in \bag''(y')$, where the nodes $y$ and $y'$ are non-equal and non-adjacent in $T''$.
    Let $x,x'\in V(T)$ be such that $u\in \bag'(x)$ and $v\in \bag'(x')$.
    By \eqref{eq:bobr}, we infer that $\dist_T(x,x')>p$.
    Since $u\in \bag'(x)$, we have that $\varphi(u)\in \bag(x)$, and similarly $\varphi(v)\in \bag(x')$.
    Since $\Tt$ is a tree-partition of $H$ and every edge of $H$ has weight at least $\gamma r$, the assertion $\dist_T(x,x')>p$ implies that
    \[\dist_H(\varphi(u),\varphi(v))>p\cdot \gamma r\geq (\alpha+\beta)r.\]
    Recall now that $H$ is an $(\alpha,\beta r)$-quasi-isometry. Hence,
    \[\dist_H(\varphi(u),\varphi(v))\leq \alpha\cdot \dist_G(u,v)+\beta r.\]
    By combining the two inequalities above we conclude that $\dist_G(u,v)>r$, as required.
\end{proof}

Let us stress that in the equivalences provided by \cref{thm:equivalences}, the constants $k_1$, $\Delta_1$, $k_2$, $k_3$, $k_4$, $\Delta_4$, $k_5$,$\Delta_5$, $\alpha$, $\beta$, $\gamma$ can be bounded in terms of each other and of $m$, but the distance parameter $r$ is not involved in those bounds. In other words, the equivalence holds for any choice of the ``scale'' $r\in \N_{>0}$. Also, we note that only implications \ref{pr:bsn}$\Rightarrow$ \ref{pr:qi} and \ref{pr:qia}$\Rightarrow$ \ref{pr:tpw} make use of the assumption that the doubling dimension is bounded.
    
    %and let $I$ be an inclusion-wise maximal distance-$\ell$ independent set in $T$. (We consider here the edges of $T$ to have unit weights.) For each $z\in I$, we let $D_z\subseteq V(T)$ be the set of those nodes $y$ of $T$ for which $x$ is the closest (in $T$) node belonging to $I$; ties are broken arbitrarily. Note that $\{D_z\colon z\in I\}$ is a partition of $V(T)$ and each $D_z$ induces a connected subtree of $T$ of radius at most $\ell$. Since $T$ has maximum degree at most $\Delta_5'$, we have
    %\[|D_z|\leq 1+\Delta_5'+(\Delta_5')^2+\ldots+(\Delta_5')^\ell\eqqcolon M.\]
    %Let now $T''$ be the tree obtained from $T$ by contracting $D_z$ onto $z$, for each $z\in I$. We endow $T''$ with a bag function $\bag''$ defined as follows:
    %\[\bag''(z)\coloneqq \bigcup_{x\in D_z} \bag'(x).\]
    %Since $\bag'(x)$ is $(k_5',r)$-coverable for each $x\in V(T)$, it follows that $\bag''(z)$ is $(k_1,r)$-coverable, where $k_1\coloneqq M\cdot k_5'$. Further, since $T$ has maximum degree at most $\Delta_5'$, $T''$ has maximum degree at most $\Delta_1\coloneqq \Delta_5'\cdot M$. It remains to show that $\Tt''\coloneqq (T'',\bag'')$ is a tree-partition of $G$ of spread $r$.
    
    
    %Thus each $\beta_x$ can be covered with $k_2 2^{m\log\alpha\beta}$ balls of radius~$r$. 
    %However, $\T'$ is not necessarily a dist-$r$ tree partition of $(X, d)$, as there may be two points which are close in $X$ but belong to distant bags of $T$. To address this, we choose an inclusion-maximal dist-$r(\alpha + \beta)$ independent set $I \subseteq V(T)$ and fix an arbitrary order on it. Then  we take a Voronoi diagram $D_V(V(T), I)$. Let us now consider $\T''=(T'', \{\beta_x\}_{x\in V(T'')})$, where:
    %\[V(T'')= I, \quad E(T'')=\{u_1u_2\mid \exists_{v_1\in R_V(u_1)}\exists_{v_2\in R_V(u_2)}\ v_1v_2\in E(T)\},\]
    %and for every $x\in V(T'')$ we have $\beta_x=\bigcup_{v\in R_V(x)} \beta_v$. 
    %Now each bag $\beta_x$ can be covered by \[k_2 2^{m\log\alpha\beta} \cdot \Delta_2^{r(\alpha+\beta)+1}:=k_3\] balls of radius $r$.

    %For this, consider any two vertices $u,v\in V(G)$ such that $\dist_G(u,v) \leq r$ and suppose for a contradiction that $\dist_{T''}(z,w)>1$, where $z,w\in V(T'')$ are such that $u\in \bag''(z)$ and $v\in \bag''(w)$. By construction it follows that if $x,y\in V(T)$ are such that $u\in \bag'(x)$ and $v\in \bag'(y)$, then $\dist_T(x,y)$
    
    %Suppose that $\varphi(u_1)$ and $\varphi(u_2)$ do not belong to the same node or to two neighboring nodes of $T''$.
    %Thus nodes of $T$ containing $\varphi(u_1)$ and $\varphi(u_2)$ must be at distance greater than $r(\alpha +\beta)$ in $T$ and therefore we have $d_H(\varphi(u_1), \varphi(u_2))>r(\alpha+\beta)$. % By the definition of the isometry we get
    %\[d(u_1,u_2) \geq \tfrac1\alpha(d_H(\varphi(u_1), \varphi(u_2))-\beta r) > \tfrac1\alpha\left(r(\alpha+\beta)-\beta r\right) = r,\] which contradicts the assumption $d(u_1,u_2) \leq r$. 

    %Therefore the constructed structure is indeed a dist-$r$ tree partition of $(X,d)$ of width $k_3$.

% \begin{theorem}
%    Let $(X,d)$ be a metric space, $G$ a graph of maximal degree bounded by $\Delta$  and $d_G$ the shortest path metric on $G$. Suppose that treewidth of $G$ is equal $t$ and there exist $(\alpha, \beta)$-quasi-isometry $\varphi\colon X\to G$. Then there exists dist-$r$ tree partition decomposition  of $(X,d)$ of width TODO.
% \end{theorem}
% \begin{proof}
   % Let $\mathcal{T}=(T, \{\beta_x\}_{x\in V(T)}\}$ be a tree-partition of $G$ of width TODO and degree TODO. Let us split $T$ into subtrees, each of radius $r$, and contract each of them into a single vertex, obtaining a new tree $T'$. 
% \end{proof}

%tutaj napisać że skupialiśmy się na tree partittion width, ale możemy przejść także na język treewidth (patrzy pierwszy dowód, że skoro mamy tree partition to i mamy treewidth a i że z tego wychodzą naturalnie separatory
%In the Theorem~\ref{thm:equivalencies} we focused on metric spaces with bounded doubling dimension and described dependencies based on the notion of tree-partition in graphs and its metric equivalent. In the following theorem we focus on treewidth.

%\begin{theorem}\label{thm:treewidth}
%    Let $(X,d)$ be a metric space with doubling dimension bounded by $m$ and let $r$ be a real number. Then the following conditions are equivalent:
%    \begin{enumerate}
%        \item there exists $k_1$ such that for every inclusion-maximal dist-$r$ independent set $I$ in $(X,d)$, the graph $H$ derived from the tuple $(X,d,r,I)$ has treewidth $k_1$
%        \item there exists $k_2$ such that there exists a tree decomposition of $(X,d)$ of $r$-width equal $k_2$
%        \item there exists $k_3$ such that $(X,d)$ has dist-$r$ separator number equal $k_3$
%    \end{enumerate}
%\end{theorem}

%\begin{proof} 
%    $\mathbb{3\Rightarrow 1}$ It follows drom the Lemma~\ref{lem:bdtw}.
    
%    $\mathbf{1\Rightarrow 2}$ Let $\T=(T,\{\beta_x\}_{x\in V(T)})$ be a tree decomposition of graph $H$. 
%    Let us fix any order on $V(H)$ and consider a Voronoi diagram on $(X,V(H))$. 
 %   Then let us consider
%    $\T'=(T,\{\beta_x'\}_{x\in V(T)})$, where $\beta_x'=\bigcup_{v\in\beta_x}R_V(v)$ for every $x\in V(T)$.
%    We show that $\T'$ is a tree decomposition of $(X,d)$. Suppose otherwise and let $A$ be a connected subset of $X$. Then the set of Voronoi regions which intersect $A$ is connected and therefore centers of these regions form a connected subgraph in graph $H$. Thus vertices of $T$ whose bags $\beta_x$ contain these centers also form a connected subtree by the definition of a tree decomposition of a graph. Hence, in the decomposition $\T'$, vertices whose bags intersect $A$ non-empty also form a connected subtree.

%    For every $v\in V(H)$ we have $R_V(v)\subseteq \B_{(X,d)}(v,r)$, as $I=V(H)$ is an inclusion-maximal dist-$r$ independent set. Thus each bag $\beta_x'$ for $x\in V(T)$ can be covered with $k_1$ balls of radius $r$ centered in~$V(H)\cap \beta_x'$. 

%    $\mathbf{2\Rightarrow 3}$ The proof of this implication follows exactly the proof of Theorem~\ref{thm:equivalencies} of the implication ${3\Rightarrow 4}$. We omit here detailed proofs for brevity.
%    \Jadwiga{Do we want details here?}
%\end{proof}
\section{General case}

In this section we prove \cref{thm:tdecomp-simple,thm:tdecomp}. In both cases, we will construct a suitable tree decomposition explicitly, using a recursive procedure similar to the one used in the classic algorithms for constructing tree decompositions of graphs, see e.g.~\cite[Section~7.6]{platypus}. \cref{thm:tdecomp-simple} is very simple: we just iteratively break the graph using balanced separators, accumulating them on the way throughout $\log n$ levels of recursion. The proof of \cref{thm:tdecomp} is much more intricate: the recursion keeps track of a separator that can be covered only by a constant number of balls, but the radii of those balls will grow (very slowly) during the recursion.

To facilitate the description of our recursive procedures, we need the following definition of a partial tree decomposition that encapsulates the task of decomposing a subgraph of the given graph. Let $G$ be a graph, $S$ be a subset of vertices of $G$, and $U$ be the vertex set of a connected component of $G - S$. Then a \emph{partial tree decomposition} of $(S, U)$ is a tree decomposition $\Tt$ of $G[S\cup U]$ with the following additional property: there is a bag of $\Tt$ that contains the whole $S$.

In the description we will use the Iverson notation: for a condition $\psi$, $[\psi]$ is equal to $1$ if $\psi$ is true, and $0$ otherwise.


\subsection{Superconstant ball count}

We first prove \cref{thm:tdecomp-simple}.
\thmdecompsimple*
%\prz[inline]{Is it worth introducing $c$ in the theorem? From the lemma we get that the number of balls is $0 + k(\lceil \log n \rceil +1) \leq k(\log n + 2)$. The current statement might wrongly suggest that some large constant is involved.}

The construction of the required tree decomposition is encapsulated in the following~lemma.

\begin{lemma}\label{lem:tdecomp_step_simple}
    Fix $r\in \N_{>0}$.
    Let $G$ be a graph with distance-$r$ balanced separator number at most $k$, let $S$ be a set of vertices of $G$, and let $U$ be the vertex set of a connected component of $G - S$. Suppose $S$ is $(\ell,r)$-coverable for some $\ell\in \N$, and $|U| \leq 2^m$ for some $m \in \N$. Then there exists a partial tree decomposition of $(S,U)$ whose every bag is $(\ell+k(m+1),r)$-coverable. 
\end{lemma}
\begin{proof}
    We proceed by induction on $|U|$. The case $|U|=0$ holds vacuously, as we assume $U$ to be a (non-empty) connected component.

    Let us move on to the induction step. Consider the following weight function $\mu\colon V(G)\to \R_{\geq 0}$: for $u\in V(G)$, we set $\mu(u) \coloneqq [u\in U]$. Since $G$ has distance-$r$ balanced separator number bounded by $k$, we can find a $(k,r)$-coverable set $Z\subseteq V(G)$ that is a balanced separator for $\mu$. This means that every connected component of $G-Z$ contains at most $|U|/2$ vertices of $U$.

    Let $\Aa$ be the family of the vertex sets of all connected components of $G[U]-Z$. Since every $A\in \Aa$ is entirely contained in one connected component of $G-Z$, we have $|A|\leq |U|/2\leq 2^{m-1}$. Noting that $S\cup Z$ is $(\ell+k,r)$-coverable, we may apply the induction assumption to the pair $(S\cup Z,A)$ for each $A\in \Aa$, thus obtaining a partial tree decomposition $\Tt_A$ of $(S\cup (Z\cap U),A)$ satisfying the following:
    \begin{itemize}[nosep]
        \item $\Tt_A$ has a node, say $x_A$, whose bag contains $S\cup (Z\cap U)$; and
        \item every bag of $\Tt_A$ is $((\ell+k)+km,r)$-coverable, hence $(\ell+k(m+1),r)$-coverable.
    \end{itemize}
    We may now combine the tree decompositions $\{\Tt_A \mid A\in \Aa\}$ into a single tree decomposition $\Tt$ by adding a new node $x$ with bag $S\cup (Z\cap U)$, and making $x$ adjacent to all the nodes $\{x_A \mid A\in \Aa\}$. It is straightforward to verify that $\Tt$ is a tree decomposition of $G[S\cup U]$. Also, the bag of $x$ contains $S$, so $\Tt$ is a partial tree decomposition of $(S,U)$. Finally, $S\cup (Z\cap U)$ is clearly $(\ell+k,r)$-coverable, hence every bag of $\Tt$ is $(\ell+k(m+1),r)$-coverable.
\end{proof}

Now, \cref{thm:tdecomp-simple} follows immediately by applying \cref{lem:tdecomp_step_simple} for $S=\emptyset$, $U=V(G)$ (assuming without loss of generality that $G$ is connected), $\ell=0$, and $m=\lceil \log n\rceil$.





\subsection{Superconstant radii}

We now proceed to the proof of \cref{thm:tdecomp}. We will need a few extra definitions. Recall that in the context of \cref{thm:tdecomp}, we assume the existence of balanced separators consisting of $k$ balls of radius~$r$, but to cover the bags of the constructed decomposition, we allow balls of varying radii. We will only use balls of radii being multiples of $r$, so call a set $\Bb$ of balls {\em{round}} if for every ball $B\in \Bb$, the radius of $B$, denoted $\rad(B)$, is a positive integer multiple of $r$. We will use the following {\em{potential}} of $\Bb$ to keep track of the growth of radii: \[\Phi(\Bb) \coloneqq \sum_{B \in \Bb} 2^{\rad(B)/r}.\]
Note that these definitions depend on the radius parameter $r\in \N_{>0}$, fixed in the context.

Also, we set
\[\Gamma\coloneqq 2000\cdot k^2 \log k.\]
This will be the bound on the number of balls needed to cover every separator of the constructed tree decomposition. Again, this definition depends on the parameter $k\in \N_{>0}$ fixed in the context.

With these definitions in place, our recursive procedure can be captured by the lemma below.

\begin{lemma}\label{lem:tdecomp_step}
    Fix $r\in \N_{>0}$.
    Let $G$ be a graph with distance-$r$ balanced separator number at most $k$, let $S$ be a set of vertices of $G$, and let $U$ be the vertex set of a connected component of $G - S$. Let $\Bb$ be a round set of at most $\Gamma$ balls whose union contains $S$. Suppose $|U| \leq 2^m$ for some $m \in \N$. Then there exists a partial tree decomposition $\Tt$ of $(S, U)$ such that each bag of $\Tt$
    can be covered with a round set of balls of size at most $\Gamma+2k$ and potential at most $\Phi(\Bb) + 4k(m+1)$.
\end{lemma}
\begin{proof}
    The proof is by induction on the size of $U$. There is no base of induction needed: the case $|U|=0$ cannot happen, for $U$ is the vertex set of a (non-empty) connected component of $G-S$.

    
    %Assume $|I \cap U| = 0$. We create a partial decomposition consisting of a single bag containing whole $U \cup S$. Let $\mathcal{B}'$ be the set of all balls of $\mathcal{B}$ with radii increased by $r$. The potential of $\mathcal{B}'$ is exactly $2^r \cdot \Phi(\mathcal{B})$. Now, we argue that this set indeed covers whole $U$. Assume otherwise, i.e., that there exists a point $x \in U$ not contained in any $B \in \mathcal{B}'$. We argue that in this case the set $I \cup \{x\}$ is a distance-$r$ independent set, contradicting that $I$ is inclusion-maximal.
    
    %Assume otherwise, i.e., that there exists a point $y \in I$ such that $\dist(x, y) \leq r$, hence there exists a curve $\gamma$ between $y$ and $x$ of length at most $r$. Since $y \not\in U$, the curve $\gamma$ intersects $S$, and hence some ball $B \in \mathcal{B}$. Let $o$ be the center of $B$ and let $z$ be any point in $\gamma \cap B$. We have
    %$$
    %\dist(x, o) \leq \dist(x, z) + \dist(z, o) \leq \dist(x, y) + \dist(z, o) \leq r + \rad(B).
    %$$
    %This means that $x$ is covered by $B$ with radius increased by $r$, which is a contradiction. Therefore $x$ is at distance more than $r$ from every point in $I$, and therefore $I$ is not inclusion-wise maximal. This proves the base case for induction.
    
    We proceed to the induction step. First, we need to massage the ball set $\Bb$ in order to achieve a certain ``sparseness'' property that will be useful later.

    Fix \[\alpha \coloneqq 2 + \ceil{\log 2k}.\] For $\ell \in \N_{>0}$, we shall say that a vertex $x \in V(G)$  is {\em{$\ell$-crowded}} with respect to $\Bb$ if there exist at least $2^{\alpha}$ balls in $\Bb$ of radius exactly $\ell r$ whose centers are at distance at most $\alpha r$ from $x$. Let $\Bb'$ denote the set obtained from $\Bb$ by replacing all such balls with a single ball of radius $r(\ell + \alpha)$ with center at~$x$. Clearly, we have
    \[|\Bb'|<|\Bb|\qquad\textrm{and}\qquad \Phi(\Bb')\leq \Phi(\Bb)-2^\alpha\cdot 2^\ell+2^{\ell+\alpha}=\Phi(\Bb).\]
    Moreover, by triangle inequality, every ball removed from $\Bb$ is entirely covered by the ball added, hence
    \[\bigcup \Bb'\supseteq \bigcup \Bb.\]
    By applying this operation repeatedly as long as there exists an $\ell$-crowded vertex for some $\ell$, we arrive at a round set of balls $\Bb'$ such that 
    \begin{itemize}[nosep]
        \item $|\Bb'| \leq |\Bb|$,
        \item $\Phi(\Bb') \leq \Phi(\Bb)$,
        \item $\bigcup \Bb'\supseteq \bigcup \Bb$, and
        \item no vertex $u\in V(G)$ is $\ell$-crowded with respect to $\Bb'$, for any $\ell\in \N_{>0}$.
    \end{itemize}
    Also, without loss of generality we may assume that no ball in $\Bb'$ is entirely contained in another ball from $\Bb'$, for the smaller ball can be just removed from $\Bb'$ without breaking any of the properties above.

    We observe the following consequence of the construction of $\Bb'$.

    \begin{claim}\label{clm:balls_no_bound}
        For every vertex $u \in V(G)$, there are at most $2\alpha \cdot 2^{\alpha}$ balls in $\Bb'$ whose centers lie at distance at most $\alpha r$ from $u$.
    \end{claim}
    \begin{claimproof}
        Let $\Bb'_u$ be the set of balls from $\Bb'$ whose centers are at distance at most $\alpha r$ from $u$. Pick any two distinct balls $B_1, B_2 \in \Bb'_u$, say with centers $o_1, o_2$  and radii $\ell_1 r, \ell_2 r$, respectively. Suppose for a moment that $\ell_1 \geq \ell_2 + 2 \alpha$. Then for every vertex $v \in B_2$, we have
        \[
            \dist(v, o_1) \leq \dist(v, o_2) + \dist(o_2, u) + \dist(u, o_1) \leq
            (\ell_2 + 2 \alpha) r \leq \ell_1 r,
        \]
        implying $B_2 \subseteq B_1$, which is a contradiction with the construction of $\Bb'$. Similarly, supposition $\ell_2\geq \ell_1+2\alpha$ also leads to a contradiction. Therefore we have $|\ell_1 - \ell_2| < 2 \alpha$. Since $B_1,B_2$ were chosen arbitrarily, we conclude that there are at most $2\alpha$ different radii among the balls of $\Bb'_u$. Since $u$ is not $\ell$-crowded with respect to $\Bb'$ for any $\ell \in \N$, every fixed radius gives rise to at most $2^{\alpha}$ balls of $\Bb'_u$. We conclude that $|\Bb'_u| \leq 2\alpha \cdot 2^{\alpha}$, as claimed.
    \end{claimproof}

    With $\Bb'$ constructed, we proceed with the proof.
    Since the balls of $\Bb'$ are not contained one in another, they have pairwise different centers. Let then $O$ be the set of centers of balls from $\Bb'$.
    We define the following weight functions: for $u\in V(G)$, we set \[\mu_U(u) \coloneqq [u\in U]\qquad\textrm{and}\qquad \mu_O(u) \coloneqq [u \in O].\]
    Since $G$ has distance-$r$ balanced separator number bounded by $k$, we may find sets  $\Dd_U$ and $\Dd_O$ of radius-$r$ balls with $|\Dd_U|,|\Dd_O|\leq k$ such that $\bigcup \Dd_U$ is a balanced separator for $\mu_U$ and $\bigcup \Dd_O$ is a balanced separator for $\mu_O$. Define
    \[\Dd\coloneqq \Dd_U\cup \Dd_O\qquad\textrm{and}\qquad Z\coloneqq \bigcup \Dd.\]
    Observe that $|\Dd|\leq 2k$ and $Z$ is a balanced separator both for $\mu_U$ and for $\mu_O$. 
    In what follows we will define a number of sets of balls, see \cref{fig:balls}.
%    \Jadwiga{Shall we move this reference after we define the sets?}
%    \mipilin{No, the reader should be made aware that looking at the figure while parsing definitions is helpful.}
    
    Let $O_\Dd$ be the set of centers of balls from $\Dd$. Further, let $\wh{\Dd}$ be the set of those balls from $\Bb'$ whose centers lie at distance at most $\alpha r$ from any vertex of~$O_\Dd$. By \cref{clm:balls_no_bound}, we have \[|\wh{\Dd}| \leq 2\alpha \cdot 2^{\alpha} \cdot |O_\Dd| \leq 4\alpha k \cdot 2^{\alpha}.\]

    %Recall that $G[U]$ is a connected component of $G-S$, hence $G[U]-Z=G[U]-(\bigcup \Dd\cup \bigcup \wh{\Dd})$ is broken into several connected components; let $\Ww$ be the family of the vertex sets of those components.
    Let $W$ be the vertex set of any connected component of $G-Z$. Since $Z$ is a balanced separator for $\mu_U$ and for $\mu_O$, we have \[|W \cap U| \leq |U|/2 \leq 2^{m - 1}\qquad\textrm{and}\qquad |W \cap O| \leq |O|/2\leq \Gamma / 2.\] Consider all balls of $\Bb' - \wh{\Dd}$ with centers outside of $W$ and let $R_W$ denote the largest radius among them. We set $R_W \coloneqq (\alpha - 1)r$ in case there are no such balls or all such balls have radius less than $(\alpha - 1)r$. Let $\Dd_W$ be the set of balls obtained from the set $\Dd$ by changing the radius of every ball to $R_W - (\alpha - 2)r$ (note that this value is at least $r$). 
    Next,
let $\Bb_W'$ be those balls from $\Bb'$ whose centers lie in $W$. We~set \[\Bb_W \coloneqq \wh{\Dd} \cup \Dd_W \cup \Bb'_W.\]

\begin{figure}[t]
    \centering
    \includegraphics[width=0.7\linewidth]{balls}
    \caption{Various sets of balls in the proof of \cref{thm:tdecomp}. Blue disks indicate balls from $\Dd$ and their centers, i.e., the set $O_{\Dd}$, are shown by blue dots.
    Balls  $\Dd_W$ are obtained by (possibly) enlarging balls from $\Dd$, while keeping the same centers; this is not shown in the picture for the sake of clarity.
    Red disks depict balls from $\Bb'$. Balls filled with diagonal lines have their centers in $W$, i.e., they belong to~$\Bb'_W$. Filled red disks indicate the set $\wh{\Dd}$, i.e., their centers are close to the vertices from $O_{\Dd}$.    
    }
    \label{fig:balls}
\end{figure}

 Note that
    \[
    |\mathcal{B}_W| \leq |\wh{\Dd}| + |\Dd_W| + |\mathcal{B}'_W| \leq
    4\alpha k \cdot 2^{\alpha} + 2k + \Gamma/2 \leq \Gamma,
    \]
    where the last inequality can be argued by substituting $\alpha=2+\ceil{\log 2k}$ and $\Gamma=2000\cdot k^2\log k$ followed by direct estimations.

Now let us bound the potential of $\Bb_W$. We consider two cases: either there exists a ball $B \in \Bb' - \wh{\Dd}$ with center outside of $W$ and radius at least $(\alpha - 1) r$, or not. If not, then $\Dd_W=\Dd$ and all the balls of $\Dd_W$ have radius $r$, hence 
    \[
    \Phi(\mathcal{B}_W) \leq \Phi(\wh{\Dd} \cup \mathcal{B}'_W) + \Phi(\Dd_W) \leq \Phi(\Bb') + |\Dd_W|\cdot 2\leq \Phi(\Bb) + 4k.
    \]
    If yes, then since $B \not\in \wh{\Dd} \cup \Bb'_W$, we have
    \[
    \Phi(\Bb_W) \leq \Phi(\wh{\Dd} \cup \Bb'_W) + \Phi(\Dd_W) \leq
    \Phi(\Bb' - \{B\}) + \Phi(\Dd_W) \leq
    \Phi(\Bb') - 2^{R_W/r} + 2k \cdot 2^{R_W/r - (\alpha - 2)}.
    \]
    We have $\alpha - 2 = \ceil{\log 2k}$, so
    \[
    \Phi(\Bb_W) \leq \Phi(\Bb') - 2^{R_W/r} + 2^{R_W/r} \cdot (2k \cdot 2^{-\log 2k}) = \Phi(\Bb')\leq \Phi(\Bb).
    \]
   

    The intuition of the remainder of the proof is as follows. Recall that $G[U]$ is a connected component of $G-S$, hence the removal of $Z$ breaks $G[U]$ into several smaller components. Each such component $G[A]$ is contained in some component of $G-Z$, say $G[W]$. We would like to apply induction for all components $G[A]$ as above, keeping $\Bb_W$ as the cover for a suitable set $S_A$ separating $A$ from the rest of the graph. However, the construction of $\Bb_W$ used only a subset of the balls from $\Bb'$, so we need to make sure that the part of $S$ contained in $S_A$ is still covered by $\Bb_W$.
    %\Jadwiga{Well, we never write that we discard something, so this statement is a bit misleading.}
    %\mipilin{Rephrased.}
    %\tara{Maybe at the end say something slightly more precise, like, ``so we need to make sure that the part of $S$ contained in $S_A$ is still covered by $B_W$.'' ?}
    We do this in the following claim.

    
    %Intuitively, we will construct a partial decomposition of $(S, U)$ by ``gluing'' decompositions obtained by induction for all connected components $W$ using the separator $Z$. The set $\Bb_W$ will be used as the cover for the root bag of such smaller decomposition, hence we must argue that no balls from outside of $\mathcal{B}_W$ will ``leak'' into $W$. More formally, we need the following claim.
    
    \begin{claim}\label{clm:no_leaks}
        Let $W$ be the vertex set of a connected component of $G-Z$. Then every ball $B \in \Bb' - \wh{\Dd}$ whose center lies outside of $W$ is disjoint from $W - \bigcup \Dd_W$.
    \end{claim}

    \begin{claimproof}
        Pick any ball $B\in \Bb'-\wh{\Dd}$, say of radius $r' \leq R_W$, and let $o$ be its center; assume $o\notin W$. Pick any vertex $x \in W$ with $\dist(x, o) \leq r'$ and let $P$ be a shortest path connecting $o$ and $x$. Since $o \not\in W$ and $x\in W$, there exists a ball $B' \in \Dd$ which intersects $P$; recall that $B'$ has radius $r$.
        Let $o'$ denote the center of $B'$ and let $z$ be any vertex on $P$ that belongs to $B'$.
        As $P$ is a shortest path, we have that $\dist(o,x) = \dist(o,z) + \dist(z,x)$, and since $z \in B'$, we have $\dist(o',z)\leq r$.
        Since $B \not\in \wh{\Dd}$, we have
        $
        \dist(o, z) \geq \dist(o, o') - \dist(o', z) \geq (\alpha - 1) r,
        $
        hence
        \begin{align*}
        \dist(x, o') \leq \ &  \dist(o', z) + \dist(z, x) = \dist(o', z) + \dist(o, x) - \dist(o,z)\\
        \leq \ &  r + r' - \dist(o, z) \leq  r + R_W - (\alpha - 1)r = R_W - (\alpha - 2)r.
        \end{align*}
        In particular, the ball of radius $R_W - (\alpha - 2)r$ with center at $o'$ both contains $x$ and belongs to $\Dd_W$. So $x\in \bigcup \Dd_W$ are the proof is complete.
    \end{claimproof}

    Let $\Aa$ comprise the vertex sets of all the connected components of $G[U]-Z$.
    For $A\in \Aa$, let $G[W_A]$ be the connected component of $G-Z$ such that $A\subseteq W_A$. We define
    \[S_A \coloneqq (Z \cap (U \cup S)) \cup (S \cap W_A).\]
    Let us first verify that that the ball set $\Bb_{W_A}$ covers $S_A$.
    
    \begin{claim}\label{cl:coverage}
        For every $A\in \Aa$, we have $S_A\subseteq \bigcup \Bb_{W_A}$.
    \end{claim}

    \begin{claimproof}
        Pick any $x \in S_A$. If $x \in Z$, then $x$ is covered by some ball in $\Dd_{W_A}\subseteq \Bb_{W_A}$ (recall that every ball of $\Dd_{W_A}$ is obtained from a ball of $\Dd$ by possibly increasing the radius). 
        %\tara{I think these $W$ should be $W_A$?} 
        Now, assume that $x \in S \cap W_A$, so in particular there exists a ball $B \in \Bb'$ that contains $x$. If the center of $B$ lies in $W_A$, then $B\in \Bb'_{W_A}\subseteq \Bb_{W_A}$ and consequently $x\in \bigcup \Bb_{W_A}$. If $B\in \wh{\Dd}$, then $x\in \bigcup \wh{\Dd}\subseteq \bigcup \Bb_{W_A}$. And if $B\in \Bb'-\wh{\Dd}$ and the center of $B$ lies outside of $W$, then by \cref{clm:no_leaks} we have $x\in \bigcup \Dd_{W_A}\subseteq \bigcup \Bb_{W_A}$.
    \end{claimproof}

    Observe that since $|W\cap U|\leq |U|/2$ for each component $W$ of $G-Z$, we also have $|A|\leq |U|/2\leq 2^{m-1}$ for each $A\in \Aa$. Also, by \cref{cl:coverage} we have that $\Bb_{W_A}$ covers $S_A$, and recall that $|\Bb_{W_A}|\leq \Gamma$ and $\Phi(\Bb_{W_A})\leq \Phi(\Bb)+4k$.
    Therefore, we may apply induction to the pair $(S_A,A)$ for each $A\in \Aa$, thus obtaining a partial tree decomposition $\Tt_A$ of $(S_A,A)$ with the following properties:
    \begin{itemize}[nosep]
        \item $\Tt_A$ has a node, say $x_A$, whose bag contains the whole $S_A$; and
        \item every bag of $\Tt_A$ can be covered by a round set of at most $\Gamma+2k$ balls with potential bounded by
        \[\Phi(\Bb_{W_A})+4km\leq \Phi(\Bb)+4k+4km=\Phi(\Bb)+4k(m+1).\]
    \end{itemize}
    We now combine the decompositions $\{\Tt_A \mid A\in \Aa\}$ into a tree decomposition $\Tt$ of $G[S\cup U]$ by creating a fresh node $x$ with bag $S\cup (Z\cap U)$ and making it adjacent to all the nodes $x_A$ for $A\in \Aa$. It is straightforward to verify that $\Tt$ is a tree decomposition of $G[S\cup U]$; we leave the verification to the reader. Since $S$ is contained in the bag of $x$, $\Tt$ is a partial  tree decomposition of $(S,U)$.

    It only remains to check whether the bag of $x$ --- namely $S\cup (Z\cap U)$ --- can be covered by a round set of at most $\Gamma+2k$ balls with potential bounded by $\Phi(\Bb)+4k(m+1)$. For this, we take $\Bb\cup \Dd$. Note that
    \[\bigcup (\Bb\cup \Dd)=\bigcup \Bb\cup \bigcup \Dd\supseteq S\cup Z\supseteq S\cup (Z\cap U).\]
    Finally, we have
    \[|\Bb\cup \Dd|\leq |\Bb|+|\Dd|\leq \Gamma+2k\quad \textrm{and}\quad \Phi(\Bb\cup \Dd)\leq \Phi(\Bb)+\Phi(\Dd)\leq \Phi(\Bb)+4k\leq \Phi(\Bb)+4k(m+1),\]
    where the pre-last inequality follows from each ball of $\Dd$ having potential $2$.
\end{proof}


    

    %to obtain a partial decomposition $(T_W, \beta^{(W)}_x)$. The root bag of this decomposition contains $S_W$ and every bag can be covered with at most $\Gamma$ balls of potential at most
    %$$
    %2^r \cdot \Phi(\mathcal{B}_W) + 2k(m - 1) \cdot 4^r \leq
    %2^r (\Phi(\mathcal{B}) + 2k \cdot 2^r) + 2k(m - 1) \cdot 4^r =
    %2^r \cdot \Phi(\mathcal{B}) + 2km \cdot 4^r.
    %$$
    %We repeat this for every component $W$ of $X - Z$ and construct the partial decomposition $(T, \beta_x)$ by creating a root bag containing $S \cup (Z \cap U)$ and appending as a children the root of trees $T_W$ for all $W$. Let $r_T$ denote the root vertex of $T$ and let $r_W$ denote the root vertex of $T_W$.

    %For this to work we must prove that such choice of $S_W, U_W, \mathcal{B}_W$ satisfies the assumptions of induction, i.e., that the following claim holds.

    
    %Now, we want to show that $(T, \beta_x)$ is indeed a valid partial decomposition. Clearly, the union of all bags is exactly $U \cup S$, hence it remains to show the following.
    
    %\begin{claim}
    %    Fix any connected $A \subseteq X$ and let $A_T := \{t\in V(T) \mid \beta_t \cap A \neq \varnothing\}$. Then, $A_T$ induces a connected subtree of $V(T)$.
    %\end{claim}

    %\begin{proof}
    %    By induction, $A_T \cap V(T_W)$ is connected for every $W$, hence it suffices to prove that either:
    %    \begin{itemize}
    %        \item $A_T \subseteq V(T_W)$ for some $W$ or
    %        \item $A_T$ contains $r_T$ and every $r_W$ such that $A_T \cap V(T_W) \neq \emptyset$.
    %    \end{itemize}
        
    %    First, consider a case where $A$ intersects the root bag of $T$. If $A$ intersects $Z \cap U$, then it also intersects the root bags of $T_W$ for all $W$. Otherwise, it must intersect $S$. Pick any $W$ and a bag $\beta_t \in V(T_W)$ such that $A$ intersects $\beta_t$. Any curve between a point from $S_W \cup U_W$ and a point from $S$ must intersect $S_W$, hence $r_W \in A_T$. This shows that $A_T$ satisfies the second condition.
        
    %    Now, consider a case where $A$ is disjoint from the root bag of $T$. Then either $A$ is disjoint from $S \cup U$ in which case $A_T = \emptyset$ or is contained in $W \cap U$ for some connected component $W$ of $X - Z$. By our construction, for every other component $W'$, all bags assigned to $V(T_{W'})$ are disjoint from $W$, hence $A_T \subseteq V(T_W)$.
    %\end{proof}
    
    %This proves that $(T, \beta_x)$ is indeed a valid partial tree decomposition satisfying lemma statement.

Now \cref{thm:tdecomp} follows from an easy application of \cref{lem:tdecomp_step}.

\thmdecomp*
\begin{proof}
    Assuming without loss of generality that $G$ is connected, we apply \cref{lem:tdecomp_step} to $S=\emptyset$ and $U=V(G)$. Thus, we obtain a tree decomposition $\Tt$ of $G$ whose every bag can be covered by a round set of $\Gamma+2k\leq 2002 k^2\log k$ balls whose potential is at most $4k \cdot (\ceil{\log n}+1)$. Let $\Bb$ be any such set and denote $R \coloneqq \max_{B \in \Bb} \rad(B)$. We have
    \[
    2^{R/r} \leq \Phi(\mathcal{B}) \leq 4k \cdot (\ceil{\log n}+1)\leq 12k\log n.
    \]
    Taking a logarithm, we get $R \leq r\cdot (\log k + \log \log n +\log 12)$, which completes the proof.
\end{proof}

\bibliographystyle{plain}
\bibliography{bibliography}

\end{document} 